\documentclass{situationist-preprint}

\title{The Value-Free Ideal Partially Rescued}
\msid{SCI-13-a}
\msfamily{SCI-13}
\msdiscipline{Philosophy of Science}
\mstype{Principal manuscript}
\mskeywords{value-free ideal; inductive risk; acceptance; problem choice; Rudner; Douglas; epistemic values}

\begin{document}
\maketitle

\begin{abstract}
Can science be value-free? Rudner argued that since no hypothesis is ever verified with certainty, the decision to accept one requires a judgment of how serious a mistake would be, and that judgment is ethical. Douglas has developed the argument into a general case against the value-free ideal. We argue that the ideal survives in part. Non-epistemic values legitimately guide which problems are investigated and how much evidence is demanded before a result is accepted, which is the management of inductive risk. They do not legitimately determine the content of accepted results. We show that if acceptance requires that the evidence make a hypothesis more probable than not, by a margin that values may set, then inquirers who differ in their values but share their evidence never accept incompatible results: values decide whether a result is accepted, not which result it is. And the evidential report that accompanies an accepted result, the degree of support and the methods that produced it, is the same for all. The value-free ideal is false of the conduct of science and true of its content.
\end{abstract}

\section{The question}

The value-free ideal holds that the internal reasoning of science, by which evidence bears on hypotheses, should be governed by epistemic considerations alone: by evidence and by epistemic values such as accuracy, simplicity and scope. Ethical, social and political values may determine which questions are asked and how results are used, but not which answers are accepted.

\citet{Rudner1953} argued that the ideal cannot be met. Since no hypothesis is ever completely verified, accepting it requires deciding that the evidence is strong enough, and how strong is strong enough depends on how serious a mistake would be: we demand more evidence that a drug lacks toxic side effects than that a batch of belt buckles is not defective. \citet{Jeffrey1956} replied that the scientist need not accept hypotheses at all; she should report probabilities and leave acceptance to those who act. Rudner had anticipated the reply: the attribution of a probability is itself a hypothesis that must be accepted. \citet{Levi1960} held that acceptance as belief, unlike acceptance as a basis for action, can be governed by epistemic values alone, and \citet{Hempel1965} accepted that the rules of acceptance must reflect the values attached to correct and incorrect decisions while confining them to rules rather than to particular cases.

\citet{Douglas2000,Douglas2009} revived the argument and extended it to the internal stages of research: the choice of significance levels, the characterization of ambiguous data and the choice of models all involve inductive risk, and where the consequences of error are foreseeable, scientists are morally responsible for weighing them. She distinguished a direct role of values, as reasons for accepting a claim, which she rejected, from an indirect role, in assessing whether the evidence is sufficient, which she defended. \citet{Steel2010} and \citet{Elliott2017} developed the case for the indirect role; \citet{Betz2013} defended the ideal by recommending that scientists report hedged results whose acceptance carries negligible inductive risk; \citet{Lacey1999} defended a version on which values legitimately shape strategies of research but not the acceptance of theories.

We argue that the ideal is partly rescued. Values legitimately guide the choice of problems and the management of inductive risk. They do not legitimately enter the content of accepted results. Douglas's distinction between the direct and the indirect role is the right one, and we show that the indirect role, properly understood, cannot affect what accepted results say.

\section{Framework}

\begin{definition}[Support]\label{d:support}
The \emph{support} of a hypothesis $h$ on evidence $E$ is the credence in $h$ that is rational on $E$, written $\Pr(h \mid E)$.
\end{definition}

We take support to be determined by the evidence, as EPI-17 argues: for any body of evidence there is a unique rational credence for each proposition. The argument below needs only that inquirers who share their evidence agree on support.

\begin{definition}[Acceptance]\label{d:accept}
An inquirer \emph{accepts} $h$ on $E$ with threshold $\tau$, where $\tfrac12 < \tau \le 1$, if $\Pr(h \mid E) \ge \tau$. To accept a result is to enter it into the record as established, available as a premise for further work and as a basis for action.
\end{definition}

\begin{principle}[Inductive risk]\label{p:risk}
The threshold $\tau$ for accepting a hypothesis may legitimately depend on the costs of error, including non-epistemic costs: the more serious the consequences of wrongly accepting $h$, the higher the threshold.
\end{principle}

\begin{principle}[Evidential content]\label{p:content}
The content of an accepted result, which hypothesis is accepted and the evidential report that accompanies it, is legitimate only if it is fixed by the evidence: the hypothesis accepted must be supported by the evidence to at least the threshold, and the report must state the support and the methods that produced the evidence.
\end{principle}

Principle~\ref{p:risk} is Rudner's and Douglas's point, which we grant. Principle~\ref{p:content} excludes the direct role: a hypothesis is not to be accepted because it would be good if it were true, or because its acceptance would serve a cause. The restriction of thresholds to values above one half expresses the minimal requirement that one does not accept what the evidence makes less probable than its negation.

\section{The argument}

\begin{proposition}[No value-driven conflict]\label{p:noconflict}
Two inquirers who share their evidence but differ in their thresholds never accept incompatible hypotheses.
\end{proposition}

\begin{proof}
Suppose one accepts $h$ with threshold $\tau_1$ and the other accepts $h'$ with threshold $\tau_2$, where $h$ and $h'$ are incompatible. Then $\Pr(h \mid E) \ge \tau_1 > \tfrac12$ and $\Pr(h' \mid E) \ge \tau_2 > \tfrac12$. Since they share evidence, they agree on support, and since $h$ and $h'$ are incompatible, $\Pr(h \mid E) + \Pr(h' \mid E) \le 1$, a contradiction.
\end{proof}

\begin{proposition}[Common report]\label{p:report}
Two inquirers who share their evidence issue the same evidential report on any hypothesis, whatever their thresholds.
\end{proposition}

\begin{proof}
The report states the support and the methods. Support is fixed by the evidence (Definition~\ref{d:support}), and the methods are part of the shared evidence.
\end{proof}

\begin{theorem}\label{t:main}
Under Principles~\ref{p:risk} and~\ref{p:content}, differences in non-epistemic values can make a difference to whether a result is accepted, but not to what an accepted result says.
\end{theorem}

\begin{proof}
By Principle~\ref{p:risk}, values may set the threshold, and a hypothesis whose support lies between two thresholds is accepted by one inquirer and not by the other. That is the whole of the difference. By Proposition~\ref{p:noconflict}, the results the two accept are never incompatible: one accepts $h$ where the other suspends judgment, never $h$ where the other accepts its negation or a rival. By Proposition~\ref{p:report}, the evidential report on $h$ is the same for both. So the content of whatever either accepts, the hypothesis together with its report, is fixed by the evidence.
\end{proof}

The choice of problems is a second legitimate role. Which questions are investigated, which hypotheses are formulated and tested, and which populations are studied, are rightly guided by social needs and by the interests of those affected \citep{Kitcher2001}. Problem choice determines which results there are, not what they say: once a question is asked, the answer accepted is fixed by Theorem~\ref{t:main}. Values thus act at the boundary of the record, deciding what enters it and when, while its content answers to evidence and remains open to the public criticism that SCI-01 identifies as the mark of science.

\section{Objections and replies}

\begin{objection}[Rudner's regress]
The evidential report is itself a hypothesis, about the support and the methods, and accepting it carries inductive risk. Values enter at that level.
\end{objection}

\begin{reply}
A report of support is not an empirical hypothesis about the world but a statement of what the evidence supports, and on the uniqueness view of EPI-17 it is fixed by the evidence. A report of methods describes what was done, and its inductive risk is that of ordinary observation, for which the threshold can be set as high as the practice of checking allows, since nothing is gained by setting it lower. The regress stops at reports that carry no value-dependent risk worth the name.
\end{reply}

\begin{objection}[Internal choices]
Douglas's case of ambiguous liver slides in studies of dioxin shows that values enter the characterization of data: whether an ambiguous slide is classified as a tumor depends on which error the pathologist is more anxious to avoid \citep{Douglas2000}. The data themselves are value-laden.
\end{objection}

\begin{reply}
Classifying an ambiguous slide is accepting a hypothesis about it with some threshold, and Theorem~\ref{t:main} applies at that level: values decide whether an ambiguous slide is counted, not what is seen on it. The evidential content is preserved by reporting the classification rule and the number of ambiguous cases, so that anyone with different values can recompute. Where the rule is not reported, values have entered the content, and that is a departure from the ideal, not a legitimate use of values.
\end{reply}

\begin{objection}[Background assumptions]
Evidence bears on hypotheses only relative to background assumptions, and these may be shaped by values \citep{Longino1990}. Values therefore enter the support relation itself.
\end{objection}

\begin{reply}
Background assumptions are themselves hypotheses, accepted on evidence or held as working assumptions to be tested. Values may legitimately draw attention to neglected assumptions, which is problem choice. What they cannot do is make an assumption accepted. The historical cases in which values improved science, such as the correction of androcentric assumptions in primatology, are cases in which values directed attention to alternatives that were then accepted on evidence.
\end{reply}

\begin{objection}[Acting on evidence]
If values may delay acceptance, they may delay action, and inquirers with different values will advise different policies on the same evidence. Science is then not value-free in any sense that matters to the public.
\end{objection}

\begin{reply}
The thesis concerns the content of accepted results, not the advice given on their basis. Policy advice depends on values in any case, and the common report of Proposition~\ref{p:report} is what allows the public to see where advice differs because values differ and where it does not. When credibility is assigned on grounds other than evidence, by contrast, the loss is epistemic, as EPI-16-a shows for testimonial injustice.
\end{reply}

\section{What the result does not establish}

The argument does not show that individual scientists are free of bias, or that the record of science is in fact value-free in content; it shows what the legitimate roles of values are. Nor does it show that the choice of concepts, such as a measure of well-being, is value-free; that choice belongs with problem choice and determines which questions are asked.

\section{Conclusion}

Values legitimately shape science at two points: in choosing which problems to investigate and in setting how much evidence is required before a result is accepted. Neither reaches the content of accepted results. Inquirers who share evidence but differ in values may differ in what they accept, but never accept incompatible results, and they issue the same report of what the evidence supports. The value-free ideal fails as an ideal of scientific conduct and holds as an ideal of scientific content.

\begin{thebibliography}{99}\small
\bibitem[Betz(2013)]{Betz2013} Betz, G. (2013). In defence of the value free ideal. \emph{European Journal for Philosophy of Science} 3: 207--220.
\bibitem[Douglas(2000)]{Douglas2000} Douglas, H. (2000). Inductive risk and values in science. \emph{Philosophy of Science} 67: 559--579.
\bibitem[Douglas(2009)]{Douglas2009} Douglas, H. (2009). \emph{Science, Policy, and the Value-Free Ideal}. University of Pittsburgh Press.
\bibitem[Elliott(2017)]{Elliott2017} Elliott, K.~C. (2017). \emph{A Tapestry of Values: An Introduction to Values in Science}. Oxford University Press.
\bibitem[Hempel(1965)]{Hempel1965} Hempel, C.~G. (1965). Science and human values. In \emph{Aspects of Scientific Explanation}. Free Press.
\bibitem[Jeffrey(1956)]{Jeffrey1956} Jeffrey, R.~C. (1956). Valuation and acceptance of scientific hypotheses. \emph{Philosophy of Science} 23: 237--246.
\bibitem[Kitcher(2001)]{Kitcher2001} Kitcher, P. (2001). \emph{Science, Truth, and Democracy}. Oxford University Press.
\bibitem[Lacey(1999)]{Lacey1999} Lacey, H. (1999). \emph{Is Science Value Free? Values and Scientific Understanding}. Routledge.
\bibitem[Levi(1960)]{Levi1960} Levi, I. (1960). Must the scientist make value judgments? \emph{Journal of Philosophy} 57: 345--357.
\bibitem[Longino(1990)]{Longino1990} Longino, H.~E. (1990). \emph{Science as Social Knowledge}. Princeton University Press.
\bibitem[Rudner(1953)]{Rudner1953} Rudner, R. (1953). The scientist \emph{qua} scientist makes value judgments. \emph{Philosophy of Science} 20: 1--6.
\bibitem[Steel(2010)]{Steel2010} Steel, D. (2010). Epistemic values and the argument from inductive risk. \emph{Philosophy of Science} 77: 14--34.
\end{thebibliography}

\end{document}
